% ============================================================
%  BFPME — RAPPORT DE SYNTHÈSE (PHASE 1 : DONNÉES SYNTHÉTIQUES)
%  Résumé des travaux, méthodes, résultats et illustrations
% ============================================================
\documentclass[a4paper,11pt]{article}

% ── Encodage et langue ───────────────────────────────────────
\usepackage[utf8]{inputenc}
\usepackage[T1]{fontenc}
\usepackage[french]{babel}
\usepackage{lmodern}
\usepackage{microtype}

% ── Mise en page ─────────────────────────────────────────────
\usepackage[top=2.5cm, bottom=2.5cm, left=2.5cm, right=2.5cm]{geometry}
\usepackage{setspace}
\onehalfspacing

% ── Couleurs ─────────────────────────────────────────────────
\usepackage[dvipsnames,svgnames,table]{xcolor}
% Tout le texte en noir : BfRed et DarkGray redéfinis en noir pur.
% (Les gris clairs ne servent qu'au remplissage des figures.)
\definecolor{BfRed}{RGB}{0, 0, 0}
\definecolor{DarkGray}{RGB}{0, 0, 0}
\definecolor{LightGray}{RGB}{235, 235, 235}
\definecolor{MidGray}{RGB}{150, 150, 150}

% ── Graphiques ───────────────────────────────────────────────
\usepackage{float}
\usepackage{tikz}
\usetikzlibrary{arrows.meta, positioning, shapes.geometric, calc, fit}

% ── Titres de section ────────────────────────────────────────
\usepackage{titlesec}
\titleformat{\section}{\normalfont\Large\bfseries\color{BfRed}}{\thesection}{0.6em}{}
\titleformat{\subsection}{\normalfont\large\bfseries\color{DarkGray}}{\thesubsection}{0.6em}{}
\titlespacing*{\section}{0pt}{16pt}{8pt}

% ── Tableaux et listes ───────────────────────────────────────
\usepackage{booktabs}
\usepackage{tabularx}
\usepackage{array}
\usepackage{enumitem}
\setlist[itemize]{leftmargin=1.6em, itemsep=2pt, topsep=3pt}
\newcolumntype{Y}{>{\raggedright\arraybackslash}X}
\newcolumntype{C}{>{\centering\arraybackslash}X}

% ── Mathématiques ────────────────────────────────────────────
\usepackage{amsmath, amssymb}

% ── En-tête / pied ───────────────────────────────────────────
\usepackage{fancyhdr}
\pagestyle{fancy}\fancyhf{}
\fancyhead[L]{\small\textit{BFPME — Synthèse phase~1}}
\fancyhead[R]{\small Données synthétiques}
\fancyfoot[C]{\thepage}
\renewcommand{\headrulewidth}{0.4pt}
\setlength{\headheight}{14.5pt}

% ── Légendes ─────────────────────────────────────────────────
\usepackage{caption}
\captionsetup{labelsep=endash, font=small, labelfont=bf, justification=centering}

\usepackage{hyperref}
\hypersetup{colorlinks=true, linkcolor=BfRed, urlcolor=BfRed, citecolor=BfRed}

% ── Barres TikZ réutilisables (mini bar-chart) ───────────────
% \metricbar{label}{valeur 0..1}{couleur}
\newcommand{\metricbar}[3]{%
  \node[anchor=east, font=\small] at (0, 0) {#1};
  \draw[fill=LightGray, draw=MidGray] (0.2,-0.22) rectangle (6.2,0.22);
  \draw[fill=#3, draw=#3] (0.2,-0.22) rectangle ({0.2+6.0*#2},0.22);
  \node[anchor=west, font=\small\bfseries] at (6.35,0) {#2};
}

\begin{document}

% ── Page de titre ────────────────────────────────────────────
\begin{titlepage}
\centering
\vspace*{2cm}
{\Huge\bfseries\color{BfRed} BFPME\par}
\vspace{0.4cm}
{\Large Plateforme d'aide à la décision pour le risque de crédit PME\par}
\vspace{1.2cm}
\rule{0.8\textwidth}{1.5pt}\par
\vspace{0.8cm}
{\LARGE\bfseries Rapport de synthèse — Phase~1\par}
\vspace{0.3cm}
{\large Modélisation sur données synthétiques\par}
\vspace{0.6cm}
{\large\itshape Travaux, méthodes, résultats et illustrations\par}
\vspace{1cm}
\rule{0.8\textwidth}{1.5pt}\par
\vfill
{\large Société d'accueil~: MPSoft\par}
{\large Projet de fin d'études\par}
\vspace{1cm}
\end{titlepage}

\tableofcontents
\newpage

% ============================================================
\section{Contexte et objectifs}
% ============================================================

La BFPME (Banque de Financement des PME) évalue chaque demande de financement au moyen d'un
système de scoring expert. Ce projet, mené chez MPSoft, vise à construire une plateforme
d'aide à la décision qui \textbf{prédit la probabilité de défaut (PD)} d'une PME, \textbf{explique}
chaque décision de façon auditable, et \textbf{simule} des scénarios de stress — tout en
préservant la supervision humaine du comité de crédit.

Cette première phase a été réalisée sur un \textbf{jeu de données synthétique}, faute d'accès
immédiat aux dossiers réels. L'objectif était de valider de bout en bout la cohérence technique
du pipeline (données $\rightarrow$ modèle $\rightarrow$ explication $\rightarrow$ agent) avant
la migration vers des données réelles, traitée en phase~2.

\begin{table}[H]
\centering
\caption{Périmètre livré en phase~1}
\renewcommand{\arraystretch}{1.3}
\begin{tabularx}{\textwidth}{@{}l Y@{}}
\toprule
\textbf{Brique} & \textbf{Livré} \\
\midrule
Données        & Générateur synthétique piloté par un facteur de risque latent ($\sim$90 variables) \\
Modélisation   & Comparaison Random Forest / XGBoost sur 3 périmètres de variables (A/B/C) \\
Explicabilité  & Contributions SHAP globales et locales (\texttt{TreeExplainer}) \\
Agent IA       & Orchestrateur \texttt{/chat/analyze} + synthèse LLM local (Ollama) \\
Tests de stress& Moteur Monte Carlo vectorisé (VaR, Expected Shortfall) \\
\bottomrule
\end{tabularx}
\end{table}

% ============================================================
\section{Méthodes}
% ============================================================

\subsection{Génération des données synthétiques}

Le jeu de données est construit en cinq étapes reproductibles (graine fixe \texttt{SEED=42})~:
\begin{enumerate}[leftmargin=1.8em]
  \item \textbf{Socle de scoring BFPME} — reproduction de la formule officielle
        (\texttt{ScoringService.cs}) pour dériver les scores et la classe de risque.
  \item \textbf{Facteur de risque latent} — variable cachée synthétisant la « vraie » qualité
        d'un dossier, combinant score normalisé, indice de risque et type de projet, plus un
        bruit gaussien.
  \item \textbf{Génération conditionnelle} — plus de 90~variables métier générées
        conditionnellement au profil latent (liquidité, garanties, comportement de paiement).
  \item \textbf{Stochasticité et écrêtage} — bruit additif, tirages de Bernoulli, tirages
        multinomiaux, écrêtage à des domaines plausibles.
  \item \textbf{Défaut synthétique} — cible tirée d'un score logistique latent :
        $\text{pd\_target}=\sigma(z)$, puis $\text{default\_flag}\sim\mathrm{Bernoulli}(\text{pd\_target})$.
\end{enumerate}

\subsection{Chaîne de prétraitement et ingénierie de variables}

Des ratios métier sont dérivés avant modélisation (\texttt{loan\_to\_revenue},
\texttt{cash\_to\_debt}, \texttt{payment\_stress\_index}, \texttt{garantie\_to\_loan}, etc.),
puis un \texttt{ColumnTransformer} applique un traitement par nature de variable
(figure~\ref{fig:pipeline}). Le jeu passe de \textbf{52} variables d'entrée à \textbf{80}
variables transformées.

\begin{figure}[H]
\centering
\begin{tikzpicture}[node distance=6mm and 8mm,
   box/.style={draw=DarkGray, thick, rounded corners=2pt, fill=LightGray,
               text width=2.5cm, minimum height=1cm, align=center, font=\scriptsize},
   io/.style={draw=BfRed, very thick, rounded corners=5pt, fill=white,
              text width=2.1cm, minimum height=1cm, align=center, font=\scriptsize\bfseries},
   fl/.style={-{Stealth[length=2mm]}, thick, DarkGray}]
  \node[io] (raw) {52 variables brutes};
  \node[box, right=of raw] (fe) {Ingénierie de variables (ratios)};
  \node[box, right=of fe] (win) {Winsorisation + imputation};
  \node[box, right=of win] (enc) {Encodage + standardisation};
  \node[io, right=of enc] (out) {80 variables};
  \draw[fl] (raw)--(fe); \draw[fl] (fe)--(win); \draw[fl] (win)--(enc); \draw[fl] (enc)--(out);
\end{tikzpicture}
\caption{Chaîne de prétraitement sérialisée dans l'artefact \texttt{joblib}}
\label{fig:pipeline}
\end{figure}

\subsection{Modélisation et protocole d'évaluation}

Trois \textbf{périmètres de variables} sont comparés (tableau~\ref{tab:options}), chacun évalué
avec deux familles ensemblistes~: \textbf{Random Forest} et \textbf{XGBoost}. Protocole commun~:
partition stratifiée 80/20 (\texttt{random\_state=42}), 8\,000 dossiers d'entraînement,
2\,000 de test, métriques AUC / précision / rappel / F1 / exactitude.

\begin{table}[H]
\centering
\caption{Trois options de sélection de variables}
\label{tab:options}
\renewcommand{\arraystretch}{1.3}
\begin{tabularx}{\textwidth}{@{}c Y l@{}}
\toprule
\textbf{Option} & \textbf{Périmètre} & \textbf{Statut} \\
\midrule
A & Toutes les variables + scores BFPME complets & Benchmark \\
B & Sans la famille de scores BFPME (\textit{data-driven}) & Benchmark \\
C & Hybride minimal : \texttt{score\_bfpme\_global} conservé, 62 \texttt{SF\_*} retirés & \textbf{Déployé} \\
\bottomrule
\end{tabularx}
\end{table}

\subsection{Explicabilité, agent et tests de stress}

\begin{itemize}
  \item \textbf{XAI} — un \texttt{TreeExplainer} calcule les contributions SHAP sur la matrice
        transformée (importance globale et décomposition locale par dossier).
  \item \textbf{Agent IA} — le point d'entrée \texttt{/chat/analyze} enchaîne sélectivement
        validation, prédiction, explication et scénario, puis délègue à un LLM local
        (\texttt{qwen2.5:7b} via Ollama) la \emph{seule} mise en récit, sous séparation stricte
        des rôles (le LLM ne produit aucune valeur numérique).
  \item \textbf{Monte Carlo} — un moteur vectorisé génère des milliers de scénarios choqués et
        reconstitue la distribution de la PD (VaR 95\,\%, Expected Shortfall).
\end{itemize}

% ============================================================
\section{Résultats}
% ============================================================

\subsection{Performances du modèle déployé}

Le modèle retenu est le \textbf{XGBoost de l'option~C}. Le tableau~\ref{tab:res} compare les deux
familles sur ce périmètre~; les performances sont proches, XGBoost offrant un léger avantage en
rappel et Random Forest en AUC.

\begin{table}[H]
\centering
\caption{Résultats sur l'échantillon de test (option~C, \texttt{final\_dataset})}
\label{tab:res}
\renewcommand{\arraystretch}{1.3}
\begin{tabularx}{\textwidth}{@{}l C C C C C@{}}
\toprule
\textbf{Modèle} & \textbf{AUC} & \textbf{Exactitude} & \textbf{Précision} & \textbf{Rappel} & \textbf{F1} \\
\midrule
XGBoost \textbf{(déployé)} & 0{,}889 & 0{,}855 & 0{,}868 & \textbf{0{,}875} & 0{,}872 \\
Random Forest              & \textbf{0{,}891} & 0{,}856 & 0{,}871 & 0{,}874 & 0{,}873 \\
\bottomrule
\end{tabularx}
\end{table}

\begin{figure}[H]
\centering
\begin{tikzpicture}
  \begin{scope}
    \node[anchor=east, font=\small] at (0,0) {AUC};
    \draw[fill=LightGray, draw=MidGray] (0.2,-0.2) rectangle (6.2,0.2);
    \draw[fill=BfRed, draw=BfRed] (0.2,-0.2) rectangle ({0.2+6.0*0.889},0.2);
    \node[anchor=west, font=\small\bfseries] at (6.35,0) {0{,}889};
  \end{scope}
  \begin{scope}[yshift=-0.7cm]
    \node[anchor=east, font=\small] at (0,0) {Précision};
    \draw[fill=LightGray, draw=MidGray] (0.2,-0.2) rectangle (6.2,0.2);
    \draw[fill=BfRed, draw=BfRed] (0.2,-0.2) rectangle ({0.2+6.0*0.868},0.2);
    \node[anchor=west, font=\small\bfseries] at (6.35,0) {0{,}868};
  \end{scope}
  \begin{scope}[yshift=-1.4cm]
    \node[anchor=east, font=\small] at (0,0) {Rappel};
    \draw[fill=LightGray, draw=MidGray] (0.2,-0.2) rectangle (6.2,0.2);
    \draw[fill=BfRed, draw=BfRed] (0.2,-0.2) rectangle ({0.2+6.0*0.875},0.2);
    \node[anchor=west, font=\small\bfseries] at (6.35,0) {0{,}875};
  \end{scope}
  \begin{scope}[yshift=-2.1cm]
    \node[anchor=east, font=\small] at (0,0) {F1};
    \draw[fill=LightGray, draw=MidGray] (0.2,-0.2) rectangle (6.2,0.2);
    \draw[fill=BfRed, draw=BfRed] (0.2,-0.2) rectangle ({0.2+6.0*0.872},0.2);
    \node[anchor=west, font=\small\bfseries] at (6.35,0) {0{,}872};
  \end{scope}
\end{tikzpicture}
\caption{Métriques du XGBoost déployé (option~C)}
\label{fig:barres}
\end{figure}

\subsection{Réalisme : de la v1 à la v2}

La première version du générateur produisait des performances quasi parfaites (AUC $\approx$
0{,}99), attendues sur des données trop cohérentes. Une seconde version a \emph{injecté du
réalisme}~: davantage de bruit, réduction du poids des variables dominantes, cas contradictoires,
et surtout un \textbf{découpage temporel} (train = dossiers anciens, test = dossiers récents)
introduisant une dérive de distribution. Les performances retombent alors à des niveaux
crédibles (tableau~\ref{tab:realisme}).

\begin{table}[H]
\centering
\caption{Effet de l'injection de réalisme (option~A, découpage temporel v2)}
\label{tab:realisme}
\renewcommand{\arraystretch}{1.3}
\begin{tabularx}{\textwidth}{@{}l C C C@{}}
\toprule
\textbf{Version} & \textbf{AUC (RF)} & \textbf{AUC (XGB)} & \textbf{Lecture} \\
\midrule
v1 (split aléatoire)  & $\approx$ 0{,}99 & $\approx$ 0{,}99 & Cohérence technique du pipeline \\
v2 (split temporel)   & 0{,}820 & 0{,}802 & Robustesse sous hypothèses réalistes \\
\bottomrule
\end{tabularx}
\end{table}

\noindent\textbf{Message clé~:} des métriques plus basses ne signifient pas un projet dégradé —
elles traduisent un cadre d'évaluation plus honnête et plus défendable avant la migration vers
les données réelles.

% ============================================================
\section{Conclusion et transition vers la phase~2}
% ============================================================

La phase~1 valide une chaîne complète et cohérente~: génération de données réalistes,
modélisation comparée (RF / XGBoost sur trois périmètres), explicabilité SHAP, agent IA
orchestrateur avec LLM local, et tests de stress Monte Carlo. Le modèle déployé (XGBoost
option~C) atteint une AUC de \textbf{0{,}89} sur le jeu final, et de \textbf{0{,}80–0{,}82} sous
évaluation temporelle stricte.

La suite logique — objet du \textbf{rapport de phase~2 (BFPME\_2)} — est la sortie du cadre
synthétique~: confrontation à des \textbf{jeux de données externes} pour éprouver la
généralisation, puis migration vers les \textbf{données réelles de la BFPME} pour un
ré-entraînement et une validation en conditions opérationnelles.

\end{document}
